\documentclass[a4paper,12pt]{article}
\usepackage{color}
\usepackage{graphicx, gensymb}
\usepackage{amsfonts, cancel, amssymb, mathrsfs}
\usepackage{amsmath, array, esvect, esint}
\usepackage{typearea}
\usepackage[a4paper,left=2cm,right=2cm,top=2cm,bottom=3cm,bindingoffset=5mm]{geometry}
\newcommand\numberthis{\addtocounter{equation}{1}\tag{\theequation}}

\begin{document}

\title{Major Concepts of Calculus for Physics}
\author{Joachim Schwardt}
\date{\today}
\maketitle

\section{Trigonometric Identities}
For the functions $\sin(x)$, $\cos(x)$ and $\tan(x)$ we have the following identities:
About $\sin(x)$:
\begin{align*}
\sin(x) &= \frac{e^{ix}-e^{-ix}}{2i} \\
\sin(x) &= \sum_{k=0}^{\infty}\frac{(-1)^k}{(2k+1)!}x^{2k+1} \\
\sin(x) &= x\prod_{k=1}^{\infty}\left(1-\frac{x^2}{k^2\pi^2}\right) \\
\sin^2(x) &= \frac{1}{2}(1-\cos(2x)) \\
\sin(2x) &= 2\sin(x)\cos(x) \\
\sin\left(\frac{x}{2}\right) &= \sqrt{\frac{1-\cos(x)}{2}} \\
\frac{a}{\sin(\alpha)} &= \frac{b}{\sin(\beta)}=\frac{c}{\sin(\gamma)} \\
\sin(x\pm y) &= \sin(x)\cos(y)\pm\cos(x)\sin(y) \\
\sin(x+y)\sin(x-y)&=\sin^2(x)-\sin^2(y)=\cos^2(y)-\cos^2(x) \\
\sin(nx) &= \sin(x) \sum_{k=0}^{\lfloor\frac{n-1}{2}\rfloor} (-1)^k {n-k-1 \choose k} 2^{n-2k-1} \cos^{n-2k-1}(x) \\
\sin^n(x) &= \frac{1}{2^n}\sum_{k=0}^n {n\choose k}\cos\left( (n-2k)(x-\frac{\pi}{2}) \right) \\
\sin(x)+\sin(y) &= 2\sin\left(\frac{x+y}{2}\right)\cos\left(\frac{x-y}{2}\right) \\
\sin(x)\sin(y) &= \frac{1}{2}\left(\cos(x-y)-\cos(x+y)\right)
\end{align*}
About $\cos(x)$:
\begin{align*}
\cos(x) &= \frac{e^{ix}+e^{-ix}}{2} \\
\cos(x) &= \sum_{k=0}^{\infty}\frac{(-1)^k}{(2k)!}x^{2k} \\
\cos(x) &= \prod_{k=1}^\infty  \left(1-\frac{4x^2}{(2k-1)^2\pi^2} \right) \\
\cos^2(x) &= \frac{1}{2}(1+\cos(2x)) \\
\cos(2x) &= \cos^2(x) - \sin^2(x)\\
\cos\left(\frac{x}{2}\right) &= \sqrt{\frac{1+\cos(x)}{2}} \\
a^2 &= b^2+c^2-2bc\cos(\alpha) \\
a &= b\cos(\gamma) + c\cos(\beta) \\
\cos(x\pm y) &= \cos(x)\cos(y)\mp\sin(x)\sin(y) \\
\cos(x+y)\cos(x-y)&=\cos^2(x)-\sin^2(y)=\cos^2(y)-\sin^2(x) \\
\cos(nx) &= \sum_{k=0}^{\lfloor\frac{n}{2}\rfloor}(-1)^k{n \choose 2k}\sin^{2k}(x)\cos^{n-2k}(x) \\
\cos^n(x)&= \frac{1}{2^n}\sum_{k=0}^n {n\choose k}\cos\left( (n-2k)x \right) \\
\cos(x)+\cos(y) &= 2\cos\left(\frac{x+y}{2}\right)\cos\left(\frac{x-y}{2}\right) \\
\cos(x)\cos(y) &= \frac{1}{2}\left(\cos(x-y)+\cos(x+y)\right)
\end{align*}
About $\tan(x)$:
\begin{align*}
\tan(x) &= \frac{e^{ix}-e^{-ix}}{ie^{ix}+ie^{-ix}} \\
\tan(x) &= \prod_{k=-\infty}^\infty \left(\frac{x+k\pi}{x+k\pi+\frac\pi2} \right) \\
\tan(x\pm y)&=\frac{\tan(x)\pm\tan(y)}{1\mp\tan(x)\tan(y)} \\
\tan(x\pm y)&=\frac{\sin(x\pm y)}{\cos(x)\cos(y)}
\end{align*}
Trigonometric Pythagoras:
\begin{align*}
\sin^2(x)+\cos^2(x)=1 \\
\tan^2(x)+1=\sec^2(x) \\
1+\cot^2(x)=\csc^2(x)
\end{align*}
\section{Laplace Transformation}
We define the Laplace transform of a function as:
\begin{align*}
\mathscr{L}\{f(t)\}(s)&=\int_0^{\infty}f(t)e^{-st}\,dt
\end{align*}
\section{Fourier Transformation}
We define the Fourier transform of a function as:
\begin{align*}
\mathcal{F}[f(t)](\omega)&=\frac{1}{\sqrt{2\pi}}\int_{-\infty}^{\infty}f(t)e^{-i\omega t}\,dt =f(\omega)& \\
\mathcal{F}^{-1}[f(\omega)] &=\frac{1}{\sqrt{2\pi}}\int_{-\infty}^{\infty}f(\omega)e^{i\omega t}\,d\omega =f(t)&
\end{align*}

\end{document}